% ==================================================================
% id:       comb-003
% topic:    קומבינטוריקה
% subtopic: הסברים קומבינטוריים לזהויות
% solution: full
% verified: no
% ==================================================================

\begin{enumerate}
\item תנו הסבר קומבינטורי לשוויון
\[
\binom{30}{13}=\sum_{k=0}^{10}\binom{10}k\binom{20}{13-k}.
\]
(רמז: גפן צריכה לארוז 13 זוגות מכנסיים לחופשה. בארון יש לה 30 זוגות, מתוכם 10 ארוכים ו-20 קצרים. בכמה דרכים היא יכולה לעשות זאת?)
\begin{solbox}
\unverified
לגפן יש 30 זוגות מכנסיים, מתוכם 10 ארוכים ו-20 קצרים, והיא צריכה לארוז בדיוק 13. מצד אחד, יש לה $\binom{30}{13}$ דרכים לעשות זאת.
מצד שני, לפי עקרון החיבור נחלק למקרים לפי מספר המכנסיים הארוכים שארזה, $k=0,1,\dots,10$. אם ארזה $k$ ארוכים ו-$13-k$ קצרים, יש לה $\binom{10}k\binom{20}{13-k}$ דרכים לעשות זאת. לכן
\[
\binom{30}{13}=\binom{10}0\binom{20}{13}+\binom{10}1\binom{20}{12}+\dots+\binom{10}{10}\binom{20}3=\sum_{k=0}^{10}\binom{10}k\binom{20}{13-k}.
\]
(בדף השאלות הסכום נכתב עד $k=13$; זה אותו דבר, כי $\binom{10}k=0$ עבור $k>10$.)
\end{solbox}
\item חשבו על שיקול דומה, והוכיחו שלכל $n\in\N$ מתקיים
\[
\binom{2n}n=\sum_{k=0}^n\binom nk^2.
\]
(הכוונה: כמה זוגות מכנסיים יש עכשיו לגפן בארון, וכמה היא צריכה לארוז?)
\begin{solbox}
\unverified
לגפן יש עכשיו $2n$ זוגות מכנסיים, $n$ ארוכים ו-$n$ קצרים, והיא צריכה לארוז $n$ מהם. מצד אחד יש $\binom{2n}n$ דרכים.
מצד שני, נחלק למקרים לפי מספר הארוכים, $k=0,\dots,n$. אם ארזה $k$ ארוכים ו-$n-k$ קצרים, יש $\binom nk\binom n{n-k}=\binom nk^2$ דרכים (לפי הסימטריה $\binom n{n-k}=\binom nk$). לכן $\ans{\binom{2n}n=\sum_{k=0}^n\binom nk^2}$.
\end{solbox}
\end{enumerate}
